\documentclass[10pt,a4paper]{amsart}
\usepackage[margin=19mm]{geometry}
\usepackage{amsmath,amssymb,mathtools}
\usepackage[scr=rsfs,cal=euler]{mathalfa}
\usepackage{xfrac}
\usepackage[unicode,hidelinks,bookmarks=false]{hyperref}
\newcommand{\ie}{i.e.\ }
\newcommand{\cf}{cf.\ }
\newcommand{\resp}{respectively\ }
\newcommand{\Subsection}{\subsection}
\providecommand{\nobreakdash}{\nobreak\hbox{-}}
\setlength{\emergencystretch}{2em}
\multlinegap0pt
\usepackage{amssymb}
\usepackage{mathtools}
\usepackage{xcolor}
\newtheorem{lemma}{Lemma}
\newcommand{\dd}{\,d}
\title{A finite Guinand--Weil dictionary and archimedean tail order\\ for the truncated Weil quadratic form}
\author{Akiva Groskin}
\date{Source: arXiv:2607.02828v3 (CC BY 4.0)}
\begin{document}
\maketitle

\noindent\emph{Source and licence.} A finite Guinand--Weil dictionary and archimedean tail order\\ for the truncated Weil quadratic form. Version arXiv:2607.02828v3 (CC BY 4.0), distributed by arXiv under the Creative Commons Attribution 4.0 International licence, http://creativecommons.org/licenses/by/4.0/, as recorded in the arXiv licence metadata for this exact version and shown on the abstract page https://arxiv.org/abs/2607.02828v3. Full licence notice: \texttt{licenses/arxiv-2607.02828-CC-BY-4.0.txt}.\par\smallskip

\noindent\textit{Editorial scope.} Counted unit: the article's lemma lem:source-calculus (source0.tex lines 354--372). The article's complete original proof of it is delivered with it (source0.tex lines 374--405, one \texttt{proof} environment of 32 lines). No mathematics is written, repaired or re-derived: the statement and the proof are the article's own lines, in the article's own notation, and no retained line is substituted or re-flowed.  No further source-local context is needed: every cross-reference of every retained line resolves inside the two delivered spans.  Nothing was omitted from any retained span, so the package carries no omission record.  No retained line contains a citation, so the wrapper carries no bibliography and no bibliography entry.  No retained line contains a figure environment, an image-inclusion command or drawing code, so no image is delivered and the package declares no asset dependency.  Editorial formatting only, in addition: an AMS \texttt{amsart} preamble that restates the article's own declarations and macros used by the retained lines, namely the environments \texttt{lemma} on separate counters, together with the standard packages those lines need.  The retained environments print in this document as Lemma 1 in order of appearance, whereas the article prints the same items with the numbering of its own full document.  The complete native source of the article as distributed in the arXiv e-print for this exact version is bundled unchanged as \texttt{sources/204343/source0.tex}.
\begin{lemma}[Finite source calculus]\label{lem:source-calculus}
Let
\[
        \psi_{\alpha,\omega}(x)=\frac{\alpha}{\pi}\sin(2\pi\omega x).
\]
For the finite divided-difference matrix attached to this source,
\[
        \langle v,Q_{\alpha,\omega}v\rangle=\alpha K_v(\omega).
\]
Consequently, if $\mu$ is a finite signed Borel measure on $[0,1]$ and
\[
        \psi_\mu(x)=\frac1\pi\int_{[0,1]}\sin(2\pi\omega x)\dd\mu(\omega),
\]
then
\[
        \langle v,Q_\mu v\rangle
        =\int_{[0,1]}K_v(\omega)\dd\mu(\omega).
\]
\end{lemma}

\begin{proof}
For $m\ne n$,
\[
 (Q_{\alpha,\omega})_{mn}=
 \frac{\alpha(\sin(2\pi\omega m)-\sin(2\pi\omega n))}
      {\pi(m-n)},
\]
and the diagonal entry is
\[
        (Q_{\alpha,\omega})_{mm}=2\alpha\omega\cos(2\pi\omega m).
\]
On the other hand
\[
K_v(\omega)=
2\sum_{m,n=-N}^{N}u_m u_n e^{2\pi i n\omega}
 \int_0^\omega e^{2\pi i(m-n)t}\dd t .
\]
For $m\ne n$ the real part of the summand is
\[
u_m u_n
\frac{\sin(2\pi\omega m)-\sin(2\pi\omega n)}
     {\pi(m-n)},
\]
while for $m=n$ it is $2u_m^2\omega\cos(2\pi\omega m)$.  The imaginary parts
cancel in pairs under the even symmetry $u_{-k}=u_k$.  Multiplying by $\alpha$
gives the identity.  For a finite signed measure, each single-frequency entry
is continuous and bounded on $[0,1]$; on the diagonal,
\[
        \psi_\mu'(x)=2\int_{[0,1]}\omega\cos(2\pi\omega x)\dd\mu(\omega).
\]
The finite double sum may therefore be interchanged with the source integral.
\end{proof}

\end{document}
